\chapter{Methodology}
\label{ch:methodology}

\section{System Overview}
% TX: commodity router. RX: 3x ESP32-S3 N16R8 running espressif/esp-csi
% csi_recv_router (ping router at 100 Hz, CSI from ping replies).
% Host: serial 921600 baud -> parser -> Parquet sessions.
[TODO]

\section{CSI Acquisition}
% Frame format: 24 metadata columns + len int8 values, (imag, real) pairs.
% LLTF 64 subcarriers at 20 MHz; 52 usable (DC + guard nulls removed) —
% confirmed on real captures via detect_null_subcarriers. % TODO: confirm!
[TODO]

\section{Preprocessing Pipeline}
% Per receiver: usable-subcarrier selection -> Hampel filter (window 11,
% 3 sigma) -> moving-mean detrend (window 101) -> uniform resampling to
% 100 Hz -> 3 s windows, 1.5 s hop, max packet loss 10%. Receivers aligned
% at window level by host timestamps (independent clocks). Amplitude only.
[TODO]

\section{Dataset and Labeling}
% Scripted sessions with labels.json segments; 0.5 s boundary trimming.
% Subjects, environments, session counts: TODO after collection.
[TODO]

\section{Models}
% Baseline: statistical features (7 stats + 3 band energies per
% rx/subcarrier) -> SVM-RBF and random forest. Deep: CNN and CNN-LSTM
% over (n_rx, 300, 52) tensors; <2M parameters.
[TODO]

\section{Evaluation Protocol}
% THE core methodological claim: three splits (random-window for literature
% comparability — leaky by design; cross-session; leave-one-subject-out).
% Normalization statistics fit on train only. Falling recall reported
% separately. Fixed seeds; every run appends a provenance row (git SHA,
% config JSON) to experiments/results/.
[TODO]
